% Problem record op_06e9f0c7b3b62f3b. This TeX file is derived from
% database/problems_json/op_06e9f0c7b3b62f3b.json by scripts/sync-tex.mjs; edit the JSON record
% and rerun the script rather than editing this file.
\section{Multi-slot overhead of virtual channel conjugation}
\label{sec:op_06e9f0c7b3b62f3b}

\paragraph{Problem.}
What is the optimal quasiprobability overhead of implementing the complex
conjugate of an unknown quantum channel from $n$ queries?  Let
$\mathcal N:\mathcal L(A)\to\mathcal L(B)$ be an unknown channel with
$d_A:=\dim A$ and $d_B:=\dim B$, and fix orthonormal bases of $A$ and $B$.
The complex conjugate of $\mathcal N$ is the channel
\begin{equation}
  \mathcal N^{*}(X):=\overline{\mathcal N(\overline X)},
  \label{eq:06e9-conjugate}
\end{equation}
where the bar is entrywise complex conjugation in the fixed bases, so the
Choi operator of $\mathcal N^{*}$ is the entrywise conjugate of that of
$\mathcal N$.  An $n$-slot quantum comb is a physically realizable circuit
with $n$ open slots, each receiving one use of the unknown channel, whose
overall action is again a channel from $A$ to $B$; write
$\mathrm{Comb}_n$ for the set of such combs.  An $n$-slot virtual comb is
a real linear combination $\widetilde{\mathcal C}=\sum_i c_i\mathcal C_i$
with $\mathcal C_i\in\mathrm{Comb}_n$, and its base norm
\begin{equation}
  \|\widetilde{\mathcal C}\|_{\mathrm{base}}
  :=\min\Bigl\{\sum_i|c_i|:
    \widetilde{\mathcal C}=\sum_ic_i\mathcal C_i,\
    c_i\in\mathbb R,\ \mathcal C_i\in\mathrm{Comb}_n\Bigr\}
  \label{eq:06e9-base-norm}
\end{equation}
is the sampling overhead: estimating an expectation value of the output of
$\widetilde{\mathcal C}$ to additive error $\varepsilon$ by Monte Carlo
sampling of the $\mathcal C_i$ costs
$O(\|\widetilde{\mathcal C}\|_{\mathrm{base}}^{2}\varepsilon^{-2})$ runs.
Define the optimal $n$-query overhead of universal conjugation by
\begin{equation}
  g_n(d_A,d_B)
  :=\inf\Bigl\{\|\widetilde{\mathcal C}\|_{\mathrm{base}}:
    \widetilde{\mathcal C}\text{ is an $n$-slot virtual comb with }
    \widetilde{\mathcal C}(\mathcal N^{\otimes n})=\mathcal N^{*}
    \text{ for every channel }\mathcal N\Bigr\}.
  \label{eq:06e9-overhead}
\end{equation}
Since the extra slots may be discarded, $g_n\leq g_1$.  Determine
$g_n(d_A,d_B)$ in Eq.~\eqref{eq:06e9-overhead} for $n\geq2$: is
$g_n(d_A,d_B)<g_1(d_A,d_B)$ for some $n$, and what is
$\inf_{n}g_n(d_A,d_B)$?

\subsection*{Status}

Unsolved

\subsection*{Source}

The question is implicit in Zhu, Tang, Zhen, Li, Bai, and Wang, who
determine $g_1$ exactly and name multi-slot virtual protocols as future
work \sourcecite{ref:06e9-zhu-et-al}{ZTZ+26}.

\subsection*{Progress}

\begin{itemize}
  \item No completely positive supermap using any finite number of queries
  implements $\mathcal N^{*}$ of Eq.~\eqref{eq:06e9-conjugate} for every
  channel $\mathcal N$, and the same obstruction rules out a universal
  physical implementation of the adjoint $\mathcal N^{\dagger}$; the
  transpose, by contrast, admits a probabilistic single-query
  implementation.  Hence any universal conjugation must be virtual, with an
  overhead of the form Eq.~\eqref{eq:06e9-base-norm}
  \sourcecite{ref:06e9-zhu-et-al}{ZTZ+26}.

  \item A one-slot virtual comb implements complex conjugation, and its
  base norm is optimal among one-slot protocols:
  \begin{equation}
    g_1(d_A,d_B)=d_Ad_B-d_A+1.
    \label{eq:06e9-one-slot}
  \end{equation}
  The optimality proof in Eq.~\eqref{eq:06e9-one-slot} uses semidefinite
  duality for a single slot and does not extend to correlated multi-slot
  strategies \sourcecite{ref:06e9-zhu-et-al}{ZTZ+26}.

  \item The state-preparation case $d_A=1$ admits a multi-slot improvement.
For $d=d_B\geq2$ and $n\geq d-1$, Theorem 3, Eqs.~(14)--(15), of
Brzi\'c, Grinko, Studzi\'nski, and Quintino supplies CPTP maps
$C_\eta$ on $n$ copies such that
$C_\eta(\rho^{\otimes n})=\eta\rho^T+(1-\eta)I/d$ for every density
operator $\rho$, for both $a=n/[n+d(d-1)]$ and $b=-1/(d-1)$
\sourcecite{ref:06e9-brzic-et-al}{BGSQ26}.
Consequently, the virtual map
\begin{equation}
  \widetilde C=\frac{1-b}{a-b}C_a+\frac{a-1}{a-b}C_b
  \quad\text{satisfies}\quad
  \widetilde C(\rho^{\otimes n})=\rho^T.
  \label{eq:06e9-state-virtual-map}
\end{equation}
These maps are valid combs for state-preparation slots. Summing the absolute
coefficients in Eq.~\eqref{eq:06e9-state-virtual-map} gives
\begin{equation}
  1\leq g_n(1,d)\leq1+\frac{2(d-1)^2}{n+d-1},
  \qquad \inf_{n\geq1}g_n(1,d)=1.
  \label{eq:06e9-state-overhead}
\end{equation}
The lower bound in Eq.~\eqref{eq:06e9-state-overhead} follows from trace
preservation, which forces the coefficients of every exact decomposition
to sum to one. Its upper bound is strictly below $g_1(1,d)=d$ when
$n>d-1$; in particular $g_2(1,2)\leq5/3<2$.
The theorem determines the physical white-noise visibility range; this
argument supplies an upper bound on the virtual overhead, not a proof of
its exact finite-$n$ optimum.

  \item Composing virtual conjugation with the probabilistic transpose
  gives black-box access to $\mathcal N^{\dagger}$ and, for a unital
  channel, estimates expectation values of the Petz recovery map to error
  $\varepsilon$ with failure probability at most $\delta$ from
  $O(d_A^{3}d_B^{3}\varepsilon^{-2}\log(1/\delta))$ samples of the channel,
  so the value of $g_n$ directly controls the cost of such applications
  \sourcecite{ref:06e9-zhu-et-al}{ZTZ+26}.
\end{itemize}

\subsection*{References}

\begin{enumerate}[label={},leftmargin=4.8em,labelsep=0.5em]
  \footnotesize
  \item[\textup{[ZTZ+26]}]\label{ref:06e9-zhu-et-al}
  C. Zhu, Z. Tang, G. Zhen, Y. Li, G. Bai, and X. Wang, ``Simulation of
  Adjoints and Petz Recovery Maps for Unknown Quantum Channels,'' arXiv
  preprint (2026).
  \href{https://arxiv.org/abs/2602.05828}{arXiv:2602.05828}.

  \item[\textup{[BGSQ26]}]\label{ref:06e9-brzic-et-al}
  V. Brzi\'c, D. Grinko, M. Studzi\'nski, and M. T. Quintino,
“Optimal pure state cloning and transposition are complementary channels,”
arXiv preprint, version 2 (2026).
\href{https://arxiv.org/abs/2603.23628v2}{arXiv:2603.23628v2}.
\end{enumerate}

\subsection*{Comment}

State preparation is included in the question. In that subcase,
Eq.~\eqref{eq:06e9-state-overhead} proves a strict multi-slot advantage and
determines the infimum over the number of queries. The exact finite-$n$
optima for state preparation and the general-channel values with $d_A>1$
remain unresolved by these results. One-slot optimality alone does not
settle correlated multi-slot protocols. The bound on the base norm measures
variance per run; each $n$-slot run itself consumes $n$ channel queries.

\subsection*{Field}

Quantum algorithm; Quantum Resource Theory

\subsection*{Topic}

Channel simulation; Superchannels and quantum combs; Quantum error mitigation

\subsection*{ID}

\texttt{op\_06e9f0c7b3b62f3b}
